% !PRINTED 716
\weberpage{716}

If we combine formula (10) with formula \S.\,193, (15) and employ the notation introduced in \S.\,194
\[
X_1 = \sum{}^n \frac{\chi_1(n)}{n}, \quad X_2 = \sum{}^n \frac{\chi_2(n)}{n}
\]
then there results
\[
g\,h = g_1 h_1 \prod X_2.
\]
If we replace $g$ by its value (14), then for the class number it follows that
\weq{15}{h = k h_1,}
where
\weq{16}{
\begin{aligned}
k & = \frac{m\sqrt{q\varDelta_1}}{2^{\nu-1}\pi^\nu} \prod X_2 \quad \text{for odd } m \\
& = \frac{m\sqrt{\varDelta_1}}{2^{\nu-1}\pi^\nu} \prod X_2 \quad \text{for even } m,
\end{aligned}
}
and in which [by (10), (13)]:
\weq{17}{h_1 = \frac{\sqrt{\varDelta_1}}{2^{\nu-1}E} \prod X_1}
is the class number of the field $H_m$. In the latter expression the product extends over all characters $\chi_1$ of the group $C_1$ \gesperrt{with the exception of the principal character}.

That $k$ is an integer, and hence $h$ a multiple of $h_1$, will soon be shown. The numbers $k$ and $h_1$ are called the \gesperrt{first and second factor of the class number}.

In the case in which
\[
m = 2^\varkappa
\]
is a power of 2, we have by \S.\,176, (7):
\[
\varDelta_1 = 2^{(\varkappa-1)2^{\varkappa-2}-1},
\]
and if we substitute this value, then in this case we can write for the foregoing formulae:
\weq{18}{
\begin{aligned}
k & = \frac{\sqrt{2}\,2^{2^{\varkappa-3}(\varkappa-3)}\,2^\varkappa}{\pi^\nu} \prod X_2 \\
h_1 & = \frac{\sqrt{2}\,2^{2^{\varkappa-3}(\varkappa-3)}}{E} \prod X_1.
\end{aligned}
}
